%% kapa: exam=1 round=21 year=2010 period=2 form=A subject=appraisal_law subject_ko=부동산관계법규 seq=33 number=73
%% source: Q-Net 감정평가사 기출문제 (cst00309) artlSeq=5024181 file=extracted/1008/5024181/2010년 제21회 감정평가사 1차시험 기출문제/2교시.hwp
%% conversion: kordoc 4.24.51 -> tools/gen1.py (auto); review=proofread-gpt-sol; flags=none
\begin{question}{73}
측량․수로조사 및 지적에 관한 법령상 토지의 표시 변경에 관한 등기를 할 필요가 있는 경우로서 지적소관청이 관할 등기관서에 등기를 촉탁하여야 하는 사유가 아닌 것은?
\begin{choices}
  \item 지적공부에 등록된 지번이 변경된 경우
  \item 하나의 지번부여지역에 서로 다른 축척의 지적도가 있어 그 지역의 축척이 변경된 경우
  \item 바다로 된 토지가 원상으로 회복될 수 없어 등록말소 된 경우
  \item 지적소관청이 직권으로 지적공부의 등록사항 오류를 정정한 경우
  \item 등기부등본에 의하여 토지대장의 소유자를 변경한 경우
\end{choices}
\end{question}
\answer{5}
